Пусть $\xi_n$, где $n \in\N$, -- независимые случайные величины, причем
$$P(\xi_n = 1) = p \text{, } P(\xi_n = -1) = q = 1 - p$$.
Пусть $a < x < b$ -- целые числа, $S_n = x + \xi_1 + \ldots + \xi_n$, $S_0 = x$.
Суммы $S_n$ интерпретируются как суммарные выигрыши в моменты времени $n$ с начальным капиталом $x$.

Момент окончания игры $\tau = min\{n: S_n \in \{a, b\}\}$, вероятность
разорения есть $ P(S_{\tau} = a)$. Это вероятность падения капитала до критического уровня $a$ ранее достижения уровня $b$. 
Предельный уровень $b$
можно интерпретировать как предельный капитал второго игрока или
банка, если игра идет с банком (при достижении $b$ банк закрывается).

\begin{proof}
    Ясно, что $\tau$ -- марковский момент, но почему он конечен п.в. (т.е.
    является моментом остановки)? 
    
    Если $\tau(\omega) = +\infty$, то $S_n(\omega)$ никогда не попадает в $a$ и $b$, но вероятность этого нулевая. Это видно из того,
    что если $S_n(\omega) \neq a, b$ при всех $n$, то $a < S_n(\omega) < b$, ибо длина шага
    равна $1$. Можно непосредственно проверить, что вероятность постоянного нахождения в $(a, b)$ нулевая, но без всякой проверки это ясно
    из центральной предельной теоремы, ведь иначе с положительной вероятностью мы бы имели $\frac{S_n}{\sqrt{n}}\to 0$. 
    
    Итак, $\tau$ -- момент остановки.
    Поэтому в $a$ или $b$ сумма обязательно попадет, но надо знать вероятность попадания сначала в $a$.
    
    В случае $p = q = \frac{1}{2}$ последовательность $S_n$ -- мартингал относительно фильтрации порожденной 
    $\{\xi_n\}_{n>1}$. Мы хотели бы сказать,
    что $\E S_{\tau} = \E S_0 = x$, но момент $\tau$ не является ограниченным, так
    что теорема об остановке не работает. 
    
    Эта теорема применима к ограниченным моментам $\tau_k = \min (\tau, k)$, которые остаются марковскими. Для них $\E S_{\tau_k} = x$. Так как $\tau_k$ возрастает к $\tau$, то $S_{\tau_k} \to S_{\tau}$ п.в.
    
    Заметим, что $|S_{\tau_k}| \leq |a| + |b|$, так как $S_{\tau} \in \{a, b\}$, $S_n \in (a, b)$, если $n < \tau$. По теореме Лебега $\E S_{\tau_k} \to \E S_{\tau}$. Итак, $\E S_{\tau} = x$. При этом
    $\E S_{\tau} = a \cdot P(S_\tau = a) + b \cdot P(S_\tau = b)$, откуда
    \begin{equation*}
        P(S_\tau = a) = \frac{b-x}{b-a}
    \end{equation*}
    В случае $p \neq q$ мартингалом оказывается последовательность $X_n = (\frac{q}{p})^{S_n}$. В самом деле: в силу независимости $\xi_n$ имеем
    \begin{equation*}
        \E^{\F_n}X_{n+1} = X_n \E^{\F_n}(q/p)^{\xi_{n+1}} = X_n \E(q/p)^{\xi_{n+1}} = X_n
    \end{equation*}
    так как $\E(q/p)^{\xi_{n+1}} = p(q/p) + q(q/p)^{-1} = 1$. Как и выше, с учетом
оценки $|X_{min(n,\tau)}| \leq (q/p)^a + (q/p)^b$ получаем равенство $\E X_\tau = \E X_0 = (q/p)^x=(q/p)^a P(S_\tau=a)+(q/p)^b P(S_\tau = b)$, откуда находим
    \begin{equation*}
        P(S_\tau = a) = \frac{(q/p)^b - (q/p)^x}{(q/p)^b - (q/p)^a}
    \end{equation*}

\end{proof}